\textbf{Message passing and homophily.} GCN~\cite{kipf2016semi} aggregates symmetrically normalised neighbourhoods including a self loop, which acts as a low-pass filter. Oono and Suzuki show that very deep GCNs lose expressive power~\cite{oono2019graph}; we keep two layers throughout, so depth effects are out of scope.

\textbf{Heterophily-oriented designs.} Zhu et al.\ identify ego/neighbour separation, higher-order neighbourhoods and combining intermediate representations as design choices that help under heterophily~\cite{zhu2020beyond}; we reimplement only the first (and optionally 2-hop neighbourhoods) in a simplified form, so our ``H2GCN-style'' model is not the published architecture. GPR-GNN learns signed propagation weights~\cite{chien2020adaptive}, and LINKX shows that models that treat features and structure separately are strong on large non-homophilous graphs~\cite{lim2021large}; we do not run either. The widely used Geom-GCN heterophily datasets~\cite{pei2020geom} are, as Platonov et al.\ document, affected by duplicate nodes, which motivates our use of synthetic data~\cite{platonov2023critical}.

\textbf{When does a GNN lose to an MLP?} Ma et al.\ argue that homophily is not necessary for GCN to work and tie performance to how distinguishable the neighbourhood label distributions are~\cite{ma2021is}; Luan et al.\ propose post-aggregation similarity metrics to explain cases where GNNs lose to MLPs~\cite{luan2022revisiting}. Our experiment can be seen as a controlled version of this question in which the only things that change are the edges (with the features, labels and split of each seed held fixed across $h$).

\textbf{CSBM theory.} Deshpande et al.\ introduced the contextual SBM and characterised information-theoretic limits of combining graph and covariates~\cite{deshpande2018contextual}; Baranwal et al.\ analyse one graph convolution on a two-class version and show that it extends the linearly separable regime under homophily~\cite{baranwal2021graph}. We use the generative model but only measure finite-sample accuracy of trained models; we make no theoretical claim.
